%=============================================================================
\section{连接到集群}
%=============================================================================

\subsection{新旧集群：确保你连的是新集群}

Tufts 目前有两套 HPC 集群——\textbf{旧集群}正在逐步淘汰，\textbf{新集群}（升级版）是官方推荐使用的。两者的登录地址不同：

\begin{table}[h]
\centering
\begin{tabular}{lll}
\toprule
 & SSH 登录地址 & OnDemand 网址 \\
\midrule
\textbf{新集群}（推荐） & \texttt{login-prod.pax.tufts.edu} & \texttt{ondemand-prod.pax.tufts.edu} \\
旧集群（淘汰中） & \texttt{login.pax.tufts.edu} & \texttt{ondemand.pax.tufts.edu} \\
\bottomrule
\end{tabular}
\end{table}

\begin{warnbox}
\textbf{请始终使用新集群}（带 \texttt{-prod} 后缀的地址）。旧集群仍然可以登录，但 Tufts 官方强烈建议迁移到新集群，旧集群将来会被关闭。如果你之前的脚本或配置里用了 \texttt{login.pax.tufts.edu}（不带 \texttt{-prod}），请尽快更新。
\end{warnbox}

本教程所有内容均基于\textbf{新集群}。

\subsection{SSH 登录}

通过 SSH 连接到新集群的登录节点：

\begin{lstlisting}[language=bash]
ssh your_utln@login-prod.pax.tufts.edu
\end{lstlisting}

这里 \texttt{your\_utln} 是你的 Tufts 用户名（全小写）。如果你在校外，需要先连接 Tufts VPN（或使用跳板机，见后文）。

登录成功后，你会看到类似这样的提示符：

\begin{lstlisting}
[your_utln@login-p01 ~]$
\end{lstlisting}

\texttt{login-p01} 表示你当前在第 1 个登录节点上。这个编号不重要，3 个登录节点看到的文件完全一样。

\subsection{OnDemand 网页界面}

如果你不熟悉命令行，也可以通过浏览器访问：

\begin{center}
\url{https://ondemand-prod.pax.tufts.edu/}
\end{center}

OnDemand 提供图形化的文件管理器、终端、以及 Jupyter Notebook 等应用，适合入门和调试。

\subsection{进阶：免密码登录与跳板机（无需 VPN）}

每次登录都要输密码很烦。你可以配置 \textbf{SSH Key} 实现免密登录。

\subsubsection{第一步：生成密钥对（如果还没有）}

在你的\textbf{本地电脑}上运行：

\begin{lstlisting}[language=bash]
# 检查是否已有密钥
ls ~/.ssh/id_*.pub

# 如果没有，生成一对
ssh-keygen -t ed25519
# 一路回车即可（不设密码短语）
\end{lstlisting}

\subsubsection{第二步：把公钥传到集群上}

\begin{lstlisting}[language=bash]
ssh-copy-id your_utln@login-prod.pax.tufts.edu
\end{lstlisting}

输入一次密码后，以后就不需要了。验证：

\begin{lstlisting}[language=bash]
ssh your_utln@login-prod.pax.tufts.edu
# 应该直接进去，不再问密码
\end{lstlisting}

\begin{tipbox}
如果免密登录不生效，检查集群上的文件权限：
\begin{lstlisting}[language=bash]
chmod 700 ~/.ssh
chmod 600 ~/.ssh/authorized_keys
\end{lstlisting}
\end{tipbox}

\subsubsection{第三步（可选）：配置别名，简化命令}

每次输 \texttt{ssh your\_utln@login-prod.pax.tufts.edu} 太长了。在本地编辑 \texttt{\textasciitilde/.ssh/config}，添加：

\begin{lstlisting}
Host hpc
    HostName login-prod.pax.tufts.edu
    User your_utln
\end{lstlisting}

之后只需要 \texttt{ssh hpc} 即可登录，\texttt{rsync}、\texttt{scp} 也可以用这个别名：

\begin{lstlisting}[language=bash]
rsync -azP ./project/ hpc:~/project/
scp results.csv hpc:~/
\end{lstlisting}

\subsubsection{第四步（可选）：通过跳板机连接，不需要 VPN}

如果你在校外又不想用 VPN，可以通过一台\textbf{校内能访问的服务器}作为跳板（ProxyJump）。前提是你有一台可以从外网访问的校内机器（比如通过 ngrok 暴露的实验室服务器）。

在本地 \texttt{\textasciitilde/.ssh/config} 中配置：

\begin{lstlisting}
# 跳板机（你的校内服务器，通过 ngrok 等方式从外网可达）
Host my_jump_server
    HostName your_server_address
    User your_username
    Port your_port

# HPC，通过跳板机连接
Host hpc
    HostName login-prod.pax.tufts.edu
    User your_utln
    ProxyJump my_jump_server
\end{lstlisting}

\texttt{ProxyJump} 告诉 SSH：先连跳板机，再从跳板机连 HPC。对你来说就是一条命令：

\begin{lstlisting}[language=bash]
ssh hpc    # 自动经过跳板机，直达 HPC 登录节点
\end{lstlisting}

配合免密登录（对跳板机和 HPC 都做 \texttt{ssh-copy-id}），整个过程零密码、无 VPN。

\begin{lstlisting}[language=bash]
# 分别对跳板机和 HPC 配置免密（各输一次密码，之后永久免密）
ssh-copy-id my_jump_server
ssh-copy-id hpc
\end{lstlisting}

\begin{warnbox}
跳板机方案的前提是你有一台校内服务器可以从外网访问。如果你没有这样的机器，还是需要使用 Tufts VPN。
\end{warnbox}
